\subsection{\bt{Agent work and reuse behavior}{智能体工作与复用行为}}
\label{sec:sharing-eval}
\bi{\emph{Official-API study: tasks and time separation.} We use RelBench's task definitions, label-generation code, and metrics~\cite{robinson2024relbench}: post-votes (MAE), user-badge and user-engagement (AUC) on rel-stack; site-success and study-adverse (MAE), and study-outcome (AUC) on rel-trial. The official-API study uses the first three cutoffs of our previously prepared five-cutoff grid: January, April, and July 2020 for rel-stack, and January 2017--2019 for rel-trial. These are earlier temporal evaluations, not the final official test split. Each task starts a fresh agent session at each cutoff; the next period's rows are appended between cutoffs. We schedule 90 sessions: five methods, two databases, three tasks, and three cutoffs.}
{\emph{官方API实验：任务与时间划分。}我们使用RelBench的任务定义、标签生成代码和指标~\cite{robinson2024relbench}：rel-stack上的post-votes（MAE）、user-badge与user-engagement（AUC）；rel-trial上的site-success与study-adverse（MAE）、study-outcome（AUC）。官方API实验采用既有五截点网格中的前三个：rel-stack为2020年1月、4月和7月，rel-trial为2017--2019年1月。这些是较早的时间评估，并非最终官方测试划分。每项任务在每个截点启动新的智能体会话，截点之间追加下一期的记录。共安排90个会话：五种方法、两个数据库、三个任务、三个截点。}

\bi{\emph{Development and baselines.} Before formal evaluation, we compare two predeclared MAVRA configurations using only rel-amazon/user-churn validation scores at two cutoffs. The selected policy first validates an available task program, retrieves feature definitions as needed, and skips repeated extraction only for identical SQL already published and freshly audited. Its configuration is frozen before formal test scoring. The two development candidates consume 0.669\,M tokens in total; this one-time tuning cost is reported separately from the 90 formal sessions and is included in the study's API spending. The comparators are the same agent without memory, ACE~\cite{zhang2026ace}, AWM~\cite{wang2025awm}, and G-Memory~\cite{zhang2025gmemory}. ACE maintains a playbook per task using its official reflector and curator; AWM induces task workflows from successful sessions using the official prompt; G-Memory uses the official implementation with a database-wide shared memory. We adapt their published code and prompts to the common SQL tool interface and give all methods the same validation feedback. Reproduction details, revisions, and adaptations are in the artifact. Earlier gateway-based runs, including DeepPrep and FeatPilot comparisons, remain archived separately and are not pooled with this study.}
{\emph{开发与基线。}正式评估前，我们仅用rel-amazon/user-churn两个截点的验证分数比较两种预先确定的MAVRA配置。选中的策略先验证可用的任务方案，按需检索特征定义，并仅对已发布且本次重新通过审计的逐字相同SQL跳过重复提取。配置在正式测试评分前冻结。两个开发候选总计使用0.669\,M token；该一次性调优成本与90个正式会话分列，并计入本轮API支出。对照为同一智能体的无记忆版本，以及ACE~\cite{zhang2026ace}、AWM~\cite{wang2025awm}和G-Memory~\cite{zhang2025gmemory}。ACE用官方反思者与整理者维护每任务手册；AWM用官方提示词从成功会话归纳任务工作流；G-Memory使用官方实现，在同一数据库的智能体之间共享记忆。我们将其公开代码与提示词适配到共同的SQL工具接口，各方法获得相同类型的验证反馈。复现版本与适配细节见工件。较早的网关运行及其中的DeepPrep、FeatPilot对照单独归档，不与本轮合并。}

\bi{\emph{Common execution and accounting.} All methods use the official DeepSeek endpoint with the \code{deepseek-flash} alias, thinking enabled, and low reasoning effort; the provider documents this alias as \llmmodel{} on the study date. Agents share schema, join, read-only SQL, validation, and submission tools. Feature SQL returns one row per seed using only data no later than its time. Every validation fits a fixed LightGBM~\cite{ke2017lightgbm} on 50,000 sampled training seeds and performs a 5,000-seed temporal-leakage audit; only explicitly clean audits can supply a valid solution. The limits are 24 planner calls, five validations, two non-improving validations, and 600 seconds per feature SQL statement. From cutoff 2, the previous solution refit on current data supplies the validation target; after reaching it, one further improvement is allowed within three calls. Test labels do not enter agent feedback or memory updates, and final test scoring follows the sequence. Failed solutions receive the entity-table reference and retain all incurred costs. Total tokens include memory retrieval, extraction, and updates; reuse totals cover cutoffs 2--3. Runs share a server, so observed time is not an isolated speedup measure.}
{\emph{统一执行与计量。}所有方法均通过DeepSeek官方端点调用\code{deepseek-flash}，启用思考模式、低推理强度；提供方在实验日期将该别名记为\llmmodel{}。智能体共用看结构、连接、只读SQL、验证与提交工具。特征SQL为每个种子返回一行，只能使用不晚于种子时间的数据。每次验证在5万个抽样训练种子上拟合固定LightGBM~\cite{ke2017lightgbm}，并进行5000种子的时间泄漏审计；只有明确clean的审计结果可以提供有效方案。上限为24次规划模型调用、5次验证、连续2次不改进，以及每条特征SQL语句600秒。从截点2起，用上一方案在当前数据上的重新拟合分数作为验证目标；达到目标后，可在3次调用内再改进一次。测试标签不进入智能体反馈或记忆更新，最终测试评分在序列结束后执行。失败方案使用实体表参照分数并保留全部已发生成本。总token计入记忆检索、提取和更新，复用代价覆盖截点2--3。各组共享试验机，因此观测耗时不作为隔离加速比。}
\begin{figure}[t]
\centering
\ifbilingual
\includegraphics[width=\linewidth]{figures/reuse-work-zh}
\else
\includegraphics[width=\linewidth]{figures/reuse-work}
\fi
\caption{\bt{Observed agent work in the official-API study: mean tool calls over six reuse sessions per method and database (three tasks, cutoffs 2--3), including failures. Validation is the common harness tool. The artifact's official-study total tokens include internal memory-model work. This complete-policy comparison does not isolate feature composition.}{官方API实验观察到的智能体工作：每种方法、每个数据库六个复用会话（三个任务、截点2--3）的平均调用数，包含失败。验证指共同框架工具。工件中的官方实验总token包含内部记忆模型工作。完整策略对照未分离特征组合的作用。}}
\label{fig:reuse-work}
\Description{Two vertically stacked panels, one per database, compare all five methods. Stacked bars distinguish schema inspection, SQL probes, validation, memory tools, and submission with colors and hatching. The y-axis counts agent-issued tool calls per reuse session and starts at zero.}
\end{figure}

\bi{\emph{Where the work changes (Q2).} Figure~\ref{fig:reuse-work} decomposes the same official-API reuse sessions for no memory, ACE, AWM, G-Memory, and \system. On rel-stack, \system\ makes \OfStackExplore\ schema/probe calls per session versus \OfStackNoneExplore\ without memory, while retaining \OfStackValidate\ validation calls. On rel-trial these counts are \OfTrialExplore\ versus \OfTrialNoneExplore\ exploration calls and \OfTrialValidate\ validations. The published memory baselines make 14.5--21.3 exploration calls on rel-stack and 26.2--31.2 on rel-trial. The observed reduction is thus concentrated in schema inspection and SQL probing, with validation still performed. These counts include failed attempts and describe work distribution, not equal-cost operations or an isolated latency speedup. Internal memory-model work is included in the official study's total-token accounting, archived with all task scores.}
{\emph{工作量在哪里改变（Q2）。}图~\ref{fig:reuse-work}分解无记忆、ACE、AWM、G-Memory与\system\ 的同一组官方API复用会话。rel-stack上，\system\ 每个会话进行\OfStackExplore{}次结构检查或SQL探查，无记忆为\OfStackNoneExplore{}次，同时保留\OfStackValidate{}次验证。rel-trial分别为\OfTrialExplore{}与\OfTrialNoneExplore{}次探索，以及\OfTrialValidate{}次验证。已发表记忆基线在rel-stack进行14.5--21.3次探索，在rel-trial为26.2--31.2次。观察到的减少因此主要来自结构检查与SQL探查，同时仍执行验证。这些次数包含失败尝试，描述工作量分布，不假设调用代价相等，也不构成隔离的延迟加速结论；记忆内部的模型工作计入官方实验的总token，随全部任务分数一并归档。}

\bi{\emph{Quality and accounting context.} The complete official study retains all 90 sessions and their incurred costs. Relative to no memory, total tokens fall by \OfStackSaveAll\% on rel-stack and \OfTrialSaveAll\% on rel-trial; the reuse-stage reductions are \OfStackSaveReuse\% and \OfTrialSaveReuse\%. Across all three cutoffs, \OfFullQualityCount/6 tasks meet the predeclared quality margins. The study-adverse MAE is \OfAdverseRegression\% higher, outside the margin. \system\ produces nine valid rel-stack solutions and six valid rel-trial solutions; ACE, AWM, and G-Memory each produce seven on rel-trial. On site-success, only one of \system's three attempts is valid. The artifact retains every method's task scores, fallback outcomes, and memory costs. Each method has one sequence per database; its cutoffs are dependent observations.}
{\emph{质量与计量背景。}完整官方实验保留全部90个会话及已发生成本。相对无记忆，rel-stack与rel-trial的总token分别减少\OfStackSaveAll\%与\OfTrialSaveAll\%，复用阶段分别减少\OfStackSaveReuse\%与\OfTrialSaveReuse\%。全部三个截点上，\OfFullQualityCount/6项任务满足预设质量门槛；study-adverse MAE高\OfAdverseRegression\%，未达到门槛。\system\ 在rel-stack得到九个有效方案，在rel-trial得到六个；ACE、AWM与G-Memory在rel-trial各得到七个。site-success上，\system\ 的三次尝试仅一次有效。工件保留各方法全部任务分数、回退结果与记忆开销。每种方法在每个数据库只有一条序列，截点之间并非独立观测。}

\bi{\emph{Reuse trace audit.} Of twelve reuse sessions, nine first validate a returned complete task SQL and six retain it verbatim. None directly validates unchanged definition-compiled SQL. Although explicit \code{compose\_features} calls are absent, \code{find\_features} and \code{find\_program} compile internally; returning compiled SQL does not establish adoption. At the first cutoff, user-engagement directly validates eight compiled features. The temporal audit rejects three post-related columns that passed publication checks on user-badge; a rewritten SQL passes. Source-instance checks therefore do not certify every target population. These traces do not isolate the incremental quality or safety benefit of feature composition.}
{\emph{复用轨迹审计。}十二个复用会话中，九个首先验证返回的完整任务SQL，六个原样保留它；没有会话直接验证原样返回的定义编译SQL。虽然无显式\code{compose\_features}调用，\code{find\_features}和\code{find\_program}仍在内部编译；返回编译SQL不能证明采用。首截点user-engagement直接验证八个编译特征，时间审计拒绝了三个在user-badge发布检查中通过的帖子相关列，改写后的SQL通过。来源样本检查因而不能认证所有目标总体；这些轨迹未分离特征组合的额外质量或安全收益。}
